% Problem-section file following ../_template.tex.
\section{Maximum number of mutually unbiased bases}
\label{sec:problem-44}

\paragraph{Problem.}
For each integer $d\geq2$, determine the maximum number $\mu(d)$ of pairwise
mutually unbiased orthonormal bases of $\mathbb C^d$.  Two orthonormal bases
$\mathcal B_r=\{|e_i^{(r)}\rangle\}_{i=1}^{d}$ and
$\mathcal B_s=\{|e_j^{(s)}\rangle\}_{j=1}^{d}$ are mutually unbiased when
\begin{equation}
  \bigl|\langle e_i^{(r)}|e_j^{(s)}\rangle\bigr|^2=\frac1d
  \qquad\text{for every }i,j\in\{1,\ldots,d\}.
  \label{eq:p44-mutual-unbiasedness}
\end{equation}
Thus the extremal quantity defined by Eq.~\eqref{eq:p44-mutual-unbiasedness}
is
\begin{equation}
  \mu(d):=\max\left\{m:\text{there exist $m$ orthonormal bases of
  $\mathbb C^d$ that are pairwise mutually unbiased}\right\}.
  \label{eq:p44-maximum-mubs}
\end{equation}
Determine Eq.~\eqref{eq:p44-maximum-mubs} in the non-prime-power regime, in
particular for $d\in\{6,10,12,14,15\}$.

\subsection*{Status}

Unsolved

\subsection*{Source}

Kr\"uger and Werner explicitly pose determination of the maximum number of
mutually unbiased bases in arbitrary dimension; McNulty and Weigert retain
the composite-dimensional cases as open
\sourcecite{ref:p44-krueger-werner}{KW05},
\sourcecite{ref:p44-mcnulty-weigert}{MW26}.

\subsection*{Progress}

\begin{itemize}
  \item The universal dimension bound is $\mu(d)\leq d+1$, and finite-field
  constructions attain it whenever $d$ is a prime power.  Hence
  $\mu(d)=d+1$ throughout the prime-power regime
  \sourcecite{ref:p44-wootters-fields}{WF89},
  \sourcecite{ref:p44-klappenecker-roetteler}{KR04}.

  \item If $d=\prod_r p_r^{a_r}$ is the prime-power factorization, tensor
  products give $\mu(d)\geq1+\min_r p_r^{a_r}$.  Moreover, any family of $d$
  mutually unbiased bases extends to $d+1$ bases.  Consequently
  $\mu(6)\in\{3,4,5,7\}$; in particular, six bases cannot be maximal
  \sourcecite{ref:p44-klappenecker-roetteler}{KR04},
  \sourcecite{ref:p44-weiner}{Wei13}.

  \item No fourth mutually unbiased basis in $\mathbb C^6$ is known.
  Numerical searches support $\mu(6)=3$, while symmetry-reduced
  semidefinite hierarchies provide convergent certificate frameworks but have
  not resolved the unrestricted case
  \sourcecite{ref:p44-butterley-hall}{BH07},
  \sourcecite{ref:p44-gribling-polak}{GP24},
  \sourcecite{ref:p44-mcnulty-weigert}{MW26}.
\end{itemize}

\subsection*{References}

\begin{enumerate}[label={},leftmargin=4.8em,labelsep=0.5em]
  \footnotesize
  \item[\textup{[WF89]}]\label{ref:p44-wootters-fields}
  W. K. Wootters and B. D. Fields,
  ``Optimal State-Determination by Mutually Unbiased Measurements,''
  \emph{Annals of Physics} \textbf{191}, 363--381 (1989).
  \href{https://doi.org/10.1016/0003-4916(89)90322-9}%
       {doi:10.1016/0003-4916(89)90322-9}.

  \item[\textup{[KR04]}]\label{ref:p44-klappenecker-roetteler}
  A. Klappenecker and M. R\"otteler,
  ``Constructions of Mutually Unbiased Bases,'' in
  \emph{Finite Fields and Applications}, LNCS \textbf{2948}, 137--144
  (Springer, 2004).
  \href{https://doi.org/10.1007/978-3-540-24633-6_10}%
       {doi:10.1007/978-3-540-24633-6\_10};
  \href{https://arxiv.org/abs/quant-ph/0309120}{arXiv:quant-ph/0309120}.

  \item[\textup{[Wei13]}]\label{ref:p44-weiner}
  M. Weiner, ``A Gap for the Maximum Number of Mutually Unbiased Bases,''
  \emph{Proceedings of the American Mathematical Society} \textbf{141},
  1963--1969 (2013).
  \href{https://doi.org/10.1090/S0002-9939-2013-11487-5}%
       {doi:10.1090/S0002-9939-2013-11487-5};
  \href{https://arxiv.org/abs/0902.0635}{arXiv:0902.0635}.

  \item[\textup{[BH07]}]\label{ref:p44-butterley-hall}
  P. Butterley and W. Hall,
  ``Numerical Evidence for the Maximum Number of Mutually Unbiased Bases in
  Dimension Six,'' \emph{Physics Letters A} \textbf{369}, 5--8 (2007).
  \href{https://doi.org/10.1016/j.physleta.2007.04.059}%
       {doi:10.1016/j.physleta.2007.04.059};
  \href{https://arxiv.org/abs/quant-ph/0701122}{arXiv:quant-ph/0701122}.

  \item[\textup{[GP24]}]\label{ref:p44-gribling-polak}
  S. Gribling and S. Polak,
  ``Mutually Unbiased Bases: Polynomial Optimization and Symmetry,''
  \emph{Quantum} \textbf{8}, 1318 (2024).
  \href{https://doi.org/10.22331/q-2024-04-30-1318}%
       {doi:10.22331/q-2024-04-30-1318};
  \href{https://arxiv.org/abs/2111.05698}{arXiv:2111.05698}.

  \item[\textup{[MW26]}]\label{ref:p44-mcnulty-weigert}
  D. McNulty and S. Weigert,
  ``Mutually Unbiased Bases in Composite Dimensions---A Review,''
  \emph{Quantum} \textbf{10}, 2051 (2026).
  \href{https://doi.org/10.22331/q-2026-04-01-2051}%
       {doi:10.22331/q-2026-04-01-2051};
  \href{https://arxiv.org/abs/2410.23997}{arXiv:2410.23997}.

  \item[\textup{[KW05]}]\label{ref:p44-krueger-werner}
  O. Kr\"uger and R. F. Werner (eds.),
  ``Some Open Problems in Quantum Information Theory,''
  arXiv:quant-ph/0504166 (2005).
  \href{https://doi.org/10.48550/arXiv.quant-ph/0504166}%
       {doi:10.48550/arXiv.quant-ph/0504166};
  \href{https://arxiv.org/abs/quant-ph/0504166}%
       {arXiv:quant-ph/0504166}.
\end{enumerate}

\subsection*{Comment}

Determining the exact value in Eq.~\eqref{eq:p44-maximum-mubs} is stronger
than deciding whether a complete family of $d+1$ bases exists.  For example,
excluding seven bases in dimension six would not distinguish among the three
remaining possibilities $\mu(6)\in\{3,4,5\}$.  This extremal MUB problem is
distinct from the SIC existence and symmetry questions in Problems 17--19.

\subsection*{Tag}

Mutually unbiased bases; Measurement theory; Quantum state tomography;
Combinatorics; Semidefinite programming

\subsection*{ID}

\texttt{op\_c9c62042b15fcb06}
